% Generated deterministically from the audited Markdown source.
% Responsible author: Carptopus.
\documentclass[11pt]{article}
\usepackage[a4paper,margin=27mm]{geometry}
\usepackage[T1]{fontenc}
\usepackage[utf8]{inputenc}
\usepackage{lmodern}
\usepackage{microtype}
\usepackage{amsmath,amssymb,amsthm,mathtools}
\usepackage{needspace}
\usepackage{xcolor}
\usepackage[colorlinks=true,linkcolor=blue!55!black,citecolor=blue!55!black,urlcolor=blue!65!black]{hyperref}
\usepackage{xurl}
\usepackage{fancyhdr}

\newtheorem{theorem}{Theorem}[section]
\newtheorem{lemma}[theorem]{Lemma}
\newtheorem{proposition}[theorem]{Proposition}
\newtheorem{corollary}[theorem]{Corollary}

\pagestyle{fancy}
\fancyhf{}
\fancyhead[L]{Carptopus}
\fancyhead[R]{Parallel-support rank}
\fancyfoot[C]{\thepage}
\setlength{\headheight}{14pt}
\setlength{\parindent}{0pt}
\setlength{\parskip}{0.44em}
\setlength{\emergencystretch}{4em}
\urlstyle{same}

\title{Parallel-support rank and a one-sided Merino--Welsh inequality}
\author{Carptopus\\\href{mailto:carptopus@163.com}{carptopus@163.com}}
\date{4 September 2026}

\hypersetup{
  pdftitle={Parallel-support rank and a one-sided Merino-Welsh inequality},
  pdfauthor={Carptopus},
  pdfsubject={A one-sided Merino-Welsh inequality controlled by parallel-support rank},
  pdfkeywords={matroid, Tutte polynomial, Merino-Welsh conjecture, parallel class, basis activity}
}

\begin{document}
\maketitle

\begin{center}
\fcolorbox{black!35}{black!4}{\parbox{0.88\textwidth}{\centering\small
Internal candidate manuscript, version 0.1-beta. The mathematical proof and complete
manuscript have passed project-internal zero-trust audits. External mathematical review
and formal peer review are pending.}}
\end{center}

\begin{abstract}


Let \(M\) be a finite matroid without loops or coloops, and let \(P(M)\) be
the union of its nontrivial parallel classes. We prove that

\[
r_M(P(M))\ge r(M)-1
\quad\Longrightarrow\quad
T_M(0,2)\ge T_M(1,1).
\]

The inequality is strict when \(r_M(P(M))=r(M)-1\), and equality in the
whole stated class occurs exactly for direct sums of copies of \(U_{1,2}\).
The dual statement holds for nontrivial series classes. In the graphic case,
the hypothesis says that the vertex-spanning subgraph formed by parallel-edge
bundles has at most two connected components. The proof combines an exact
basis-activity formula for parallel extensions with deletion--contraction and
an exponential surplus matching in the terminal two-element cocircuit core.
The rank threshold is sharp if no further hypothesis is imposed: a binary
rank-three matroid with parallel-support rank one violates the inequality.
We also record a connected five-element counterexample at the same threshold,
showing that connectivity alone does not extend the result to corank two.

\end{abstract}

\section{Introduction}


For a finite matroid \(M\), the evaluations

\[
T_M(2,0),\qquad T_M(0,2),\qquad T_M(1,1)
\]

are central quantities in the Merino--Welsh family of inequalities. The last
evaluation is the number of bases. In the graphic case, the three evaluations
count acyclic orientations, totally cyclic orientations, and spanning trees,
respectively. The original graph conjecture asks whether a loopless,
2-connected graph satisfies \cite{merino-welsh}

\[
\max\{T_M(2,0),T_M(0,2)\}\ge T_M(1,1).
\]

The corresponding assertion for all loopless, coloopless matroids is false,
as are its additive and multiplicative variants \cite{beke-et-al}. Important positive
results are nevertheless known for such classes as series--parallel graphs
\cite{noble-royle}, split matroids \cite{ferroni-schroeter}, paving matroids \cite{chavez-lomeli-et-al}, and matroids satisfying suitable
density or cocircuit hypotheses \cite{kung}. The graph problem remains open.

This paper studies a one-sided inequality controlled by parallel structure.
Write \(P(M)\) for the union of all parallel classes of size at least two. If
\(P(M)\) spans \(M\), every rank direction has a parallel substitute, and a
direct activity comparison proves \(T_M(0,2)\ge T_M(1,1)\). The main point is
that one complete rank direction may be left outside that support. The
remaining singleton elements need not be parallel to one another, and their
number is unrestricted.

Our main theorem is the following.

\Needspace{0.18\textheight}
\begin{theorem}
\label{thm:main}

Let \(M\) be a finite matroid without loops or coloops. If

\[
r_M(P(M))\ge r(M)-1,
\]

then

\[
T_M(0,2)\ge T_M(1,1).                         \tag{1.1}
\]

If \(r_M(P(M))=r(M)-1\), then (1.1) is strict. Equality in the full class of
Theorem~\ref{thm:main} occurs exactly when

\[
M\cong\bigoplus_{i=1}^{r(M)}U_{1,2},           \tag{1.2}
\]

with the empty direct sum allowed in rank zero.

\end{theorem}

By duality, the analogous statement holds when the nontrivial series classes
span all but at most one rank direction in \(M^*\), with \(T_M(2,0)\) in
place of \(T_M(0,2)\).

The proof has two parts. First, we calculate the total contribution to
\(T_M(0,2)\) of all bases projecting to a fixed basis of the simplification.
Second, we induct on singleton parallel classes. At the final hyperplane core,
bases with zero activity contribution are grouped with bases having several
externally active parallel classes; the elementary inequality

\[
2^K\ge 1+K
\]

pays for the entire group.

The threshold is exact for a condition depending only on
\(r_M(P(M))\). Section 5 gives a rank-three binary counterexample with
\(r_M(P(M))=r(M)-2\), and also a connected counterexample of order five.
The latter shows that connectivity is not the missing hypothesis. A separate
bounded analysis reduces a possible cosimple extension to three rank-two
four-element quotient cores, but does not settle them; no corank-two or
two-disjoint-bases theorem is claimed here.

\section{Parallel classes and activities}


Let \(N=\operatorname{si}(M)\) be the simplification of a loopless matroid
\(M\). For each \(e\in E(N)\), let \(k_e\ge1\) be the size of the parallel
class of \(M\) represented by \(e\). Fix an order on \(E(N)\), order the
parallel classes accordingly, and within the class represented by \(e\) write

\[
e_1<e_2<\cdots<e_{k_e}.
\]

For a basis \(B\) of \(N\), write \(I(B)\) and \(E(B)\) for its internally
and externally active elements with respect to the chosen order.

\Needspace{0.18\textheight}
\begin{lemma}
\label{lem:activity-block}

The bases of \(M\) projecting to \(B\) contribute

\[
C_B=
2^{\sum_{f\in E(B)}k_f}
\prod_{e\in I(B)}(2^{k_e}-2)
\prod_{e\in B\setminus I(B)}(2^{k_e}-1)        \tag{2.1}
\]

to \(T_M(0,2)\). The number of such bases is

\[
W_B=\prod_{e\in B}k_e.                         \tag{2.2}
\]

\end{lemma}

\begin{proof}

A lifted basis is obtained by choosing one copy \(e_{j_e}\) from each class
represented in \(B\). Basic circuits and cocircuits behave uniformly inside
a parallel class:

\begin{itemize}
\item if \(e\in I(B)\), then the chosen copy is internally active exactly when
  \(j_e=1\);
\item if \(e\notin I(B)\), no chosen copy is internally active;
\item the \(j_e-1\) smaller unchosen copies in a selected class are externally
  active;
\item if \(f\in E(B)\), every copy in the unselected class represented by \(f\)
  is externally active.
\end{itemize}


At \((x,y)=(0,2)\), a basis with an internally active element contributes
zero, while a basis with \(j\) externally active elements contributes
\(2^j\). Summing independently over the copy chosen in each selected class
gives

\[
\sum_{j=2}^{k_e}2^{j-1}=2^{k_e}-2
\]

for an internally active class and

\[
\sum_{j=1}^{k_e}2^{j-1}=2^{k_e}-1
\]

for a non-internally-active class. The unselected externally active classes
give the leading power of two in (2.1). Formula (2.2) is immediate.
\end{proof}


\Needspace{0.18\textheight}
\begin{proposition}
\label{prop:all-heavy}

If every parallel class of a loopless matroid \(M\) has size at least two,
then

\[
T_M(0,2)\ge T_M(1,1),                          \tag{2.3}
\]

with equality exactly for direct sums of copies of \(U_{1,2}\).

\end{proposition}

\begin{proof}

For every integer \(k\ge2\),

\[
2^k-2\ge k,\qquad 2^k-1>k.
\]

Lemma~\ref{lem:activity-block} therefore gives \(C_B\ge W_B\) for every basis \(B\) of the
simplification. Summing over \(B\) proves (2.3).

For equality, every selected class in every basis must be internally active,
every class size must be two, and no basis may have an externally active
element. An ordered matroid with more than one basis has a basis that is not
entirely internally active. Thus the simplification has a unique basis. Since
the matroid is loopless, its simplification is free, and (1.2) follows. The
converse is immediate from multiplicativity of the Tutte polynomial.
\end{proof}


\section{A spanning parallel support}


We first allow singleton classes while assuming that the nontrivial classes
still span.

\Needspace{0.18\textheight}
\begin{proposition}
\label{prop:spanning-support}

Let \(M\) be a finite loopless matroid. If \(P(M)\) spans \(M\), then

\[
T_M(0,2)\ge T_M(1,1).
\]

If \(M\) has a singleton parallel class, the inequality is strict.

\end{proposition}

\begin{proof}

Induct on the number of singleton parallel classes. Proposition~\ref{prop:all-heavy} is the
base case. Choose a singleton element \(e\). Since the nontrivial parallel
classes span, \(e\) is neither a loop nor a coloop. Deletion preserves the
spanning parallel support. Contraction creates no loop because \(e\) has no
parallel mate; the images of the old nontrivial classes still span, and some
classes may merge. Both minors have fewer singleton classes, so
deletion--contraction and induction prove the inequality.

It remains to justify strictness in the only potentially marginal step. If
there is exactly one singleton \(e\) and \(M\setminus e\) is an equality case,
then \(M\setminus e\) is a direct sum of binary parallel classes. A fundamental
circuit of \(e\) meets at least two of those classes. If it meets exactly two,
contracting \(e\) merges them into a parallel class of size at least four; if
it meets at least three, the simplification of \(M/e\) has a circuit. In either
case \(M/e\) is not an equality case of Proposition~\ref{prop:all-heavy}. The contraction term
is therefore strict. Strong induction gives strictness whenever a singleton
class is present.
\end{proof}


\section{Parallel support of corank one}


We now prove Theorem~\ref{thm:main}. The spanning case is Proposition~\ref{prop:spanning-support}, so assume

\[
r_M(P(M))=r(M)-1.
\]

Put

\[
F=\operatorname{cl}_M(P(M)).
\]

Then \(F\) is a hyperplane. Since \(M\) has no coloops,
\(E(M)\setminus F\) has at least two elements.

We again induct on the number of singleton parallel classes.

\subsection{Removing singleton classes inside the hyperplane}


Suppose first that \(e\in F\) is a singleton. It is neither a loop nor a
coloop. Since \(e\in\operatorname{cl}_M(P(M))\), deletion leaves the old
nontrivial classes with rank \(r(M)-1\). No element of \(F\) becomes a coloop:
the nontrivial classes still span \(F\). Nor does an element outside \(F\)
become a coloop, because at least two such elements remain and each extends
the hyperplane.

After contracting \(e\), the old nontrivial classes have rank \(r(M)-2\) in
the rank-\(r(M)-1\) matroid \(M/e\). Contraction creates no loop because
\(e\) was a singleton, and contraction preserves the absence of coloops.
Thus both minors satisfy the hypothesis of Theorem~\ref{thm:main} and have fewer
singleton classes. The deletion minor remains in the corank-one layer, so its
inductive inequality is strict.

We may therefore assume that every element of \(F\) belongs to a nontrivial
parallel class.

\subsection{Removing surplus elements outside the hyperplane}


Every element outside \(F\) is a singleton, since a nontrivial parallel class
would be contained in \(P(M)\subseteq F\). If there are at least three such
elements, choose \(e\notin F\). Deleting \(e\) leaves at least two elements
outside \(F\), so the deletion minor remains loopless, coloopless, and in the
corank-one layer. In \(M/e\), the image of \(F\) has full rank, so Proposition
3.1 applies. Induction and deletion--contraction prove the result, and the
deletion term is strict.

It remains only to treat

\[
E(M)\setminus F=\{p,q\},                       \tag{4.1}
\]

with every class in \(F\) nontrivial.

\subsection{The terminal core}


Order all classes in \(F\) before \(p<q\). Let \(A\) be a basis of the
simplification restricted to \(F\). Both \(A\cup\{p\}\) and
\(A\cup\{q\}\) are bases.

For \(A\cup\{p\}\), the element \(p\) is internally active: the only
possible replacement outside the hyperplane is \(q\), which is later. Hence
the whole lifted basis block contributes zero to \(T_M(0,2)\). For
\(A\cup\{q\}\), the earlier element \(p\) replaces \(q\), and Lemma~\ref{lem:activity-block}
shows that the block contribution is at least its basis weight. It remains to
pay for the zero block associated with \(A\cup\{p\}\).

Take the fundamental circuit of \(q\) with respect to \(A\cup\{p\}\). It
contains \(p\), since \(A\cup\{q\}\) is a basis. It also contains a class
from \(F\), because \(p\) and \(q\) cannot be parallel. Let \(h(A)\) be the
least such class and put

\[
B(A)=A\setminus\{h(A)\}\cup\{p,q\}.            \tag{4.2}
\]

Basic exchange shows that \(B(A)\) is a basis. Relative to this basis,
\(h(A)\) is externally active, and both \(p\) and \(q\) fail to be internally
active because either may be replaced through \(h(A)\).

Fix a basis \(B\) containing \(p,q\), and let \(X_B\) be the set of classes
\(h(A)\) from all bases \(A\) mapped to \(B\). Write

\[
W_B=\prod_{e\in B\cap F}k_e,
\qquad
K_B=\sum_{h\in X_B}k_h.
\]

All classes in \(X_B\) are externally active at \(B\). Lemma~\ref{lem:activity-block} therefore
gives

\[
C_B\ge 2^{K_B}W_B.                             \tag{4.3}
\]

The basis \(B\) itself and all zero blocks mapped to it contain altogether

\[
W_B+\sum_{h\in X_B}k_hW_B=(1+K_B)W_B           \tag{4.4}
\]

bases. Since \(2^{K_B}\ge1+K_B\), (4.3) pays for (4.4). If
\(X_B\ne\varnothing\), then \(K_B\ge2\), and the inequality is strict.

A basis containing \(p,q\) with no preimage is also paid for by its own
activity block. Indeed, its \(F\)-part has size \(r(M)-2\) and can be completed
to a basis of \(F\) by an earlier class; that class can replace either \(p\)
or \(q\), so neither singleton is internally active. Lemma~\ref{lem:activity-block} gives
\(C_B\ge W_B\).

All basis blocks have now been paid for, and at least one block of the form
(4.2) occurs. The core inequality is strict. This completes the induction and
the proof of Theorem~\ref{thm:main}.
\hfill$\square$


\section{Equality, duality, and sharpness}


The corank-one case is always strict. In the spanning case, Proposition~\ref{prop:spanning-support}
and Proposition~\ref{prop:all-heavy} show that equality occurs precisely for (1.2). This proves
the equality statement in Theorem~\ref{thm:main}.

The identity

\[
T_{M^*}(x,y)=T_M(y,x)
\]

applied to Theorem~\ref{thm:main} gives the dual theorem for nontrivial series classes.

The rank threshold cannot be lowered using only the rank of \(P(M)\). Over
\(\operatorname{GF}(2)\), consider the column matroid of

\[
\begin{pmatrix}
1&0&1&0&0\\
0&1&1&0&0\\
0&0&0&1&1
\end{pmatrix}.
\]

This matroid is \(U_{2,3}\oplus U_{1,2}\). It is loopless and coloopless,
has rank three, and its only nontrivial parallel class has rank one. Hence

\[
r_M(P(M))=1=r(M)-2,
\]

but

\[
T_M(0,2)=4<6=T_M(1,1).                         \tag{5.1}
\]

Connectivity alone does not repair (5.1). Duplicate one element of
\(U_{3,4}\); equivalently, take the binary column matroid

\[
\begin{pmatrix}
1&1&0&0&1\\
0&0&1&0&1\\
0&0&0&1&1
\end{pmatrix}.
\]

It is connected, loopless, and coloopless, with parallel-support rank one.
Deletion--contraction along a member of the parallel pair gives

\[
T_M(x,y)=x^3+x^2+x+y+x^2y+xy+y^2,
\]

so

\[
T_M(0,2)=6<7=T_M(1,1).                         \tag{5.2}
\]

This second example is not cosimple: its three singleton elements form
two-element cocircuits in pairs.

For context, suppose that \(F=\operatorname{cl}_M(P(M))\) has rank
\(r(M)-2\) and that \(M\) is cosimple. The cocircuits of the rank-two quotient
\(Q=M/F\) are the cocircuits of \(M\) contained in \(E(M)\setminus F\), so
\(Q\) is cosimple. In a rank-two matroid, the complement of each parallel
class is a cocircuit. A direct case split on the number of parallel classes
therefore gives two disjoint bases of \(Q\): choose from four classes if there
are at least four; with three classes, cosimplicity supplies two classes with
two available elements; with two classes, each class has at least three
elements. A basis of \(F\) can be chosen inside \(P(M)\), and replacing each
member by a parallel mate gives a disjoint second basis of \(F\). Combining
these pairs lifts the two disjoint bases of \(Q\) to two disjoint bases of
\(M\).

This observation also gives a finite structural reduction. The induction is
made in the larger class of loopless, coloopless matroids with corank-two
parallel support for which \(Q\) contains two disjoint bases; cosimplicity is
not imposed on the minors. Fix two such bases of \(Q\). If \(|E(Q)|>4\), choose
an element \(e\) outside their union. Their lifts give two disjoint bases of
\(M\) avoiding \(e\), so \(M\setminus e\) retains full rank, has no coloop,
and remains in the same corank-two support class. Also \(e\notin F\), hence it
is a singleton and \(M/e\) is loopless; contraction preserves colooplessness.
The old parallel support has rank \(r(M)-2\) in \(M/e\), one less than the new
matroid rank, so Theorem~\ref{thm:main} controls the contraction term. Induction and
deletion--contraction therefore reduce a possible theorem for this layer to
four-element rank-two quotients. Their relevant parallel types are
\(U_{2,4}\), \((2,1,1)\), and \((2,2)\). The required quantitative activity
matching in those cores remains open in the present work. Thus (5.2) is used
only to show the sharpness of the rank-only and connectivity extensions; no
cosimple corank-two or two-disjoint-bases result is asserted.

\section{Graphic interpretation}


Let \(H\) be a connected loopless multigraph. Form a simple graph \(D(H)\)
on the same vertex set by retaining an edge \(uv\) exactly when \(H\) has at
least two parallel edges between \(u\) and \(v\). The rank of the union of
nontrivial parallel classes in the cycle matroid is

\[
|V(H)|-c(D(H)),
\]

where \(c(D(H))\) is the number of connected components, including isolated
vertices. Since \(r(M(H))=|V(H)|-1\), Theorem~\ref{thm:main} gives the following.

\Needspace{0.18\textheight}
\begin{corollary}
\label{cor:graphic}

If \(H\) is connected and bridgeless and \(D(H)\) has at most two connected
components, then the number of totally cyclic orientations of \(H\) is at
least its number of spanning trees. If \(D(H)\) has exactly two components,
the inequality is strict.

\end{corollary}

In particular, one may start with an arbitrary connected simple graph, select
a spanning forest with at most two components, and duplicate every edge of
that forest. All remaining edges may retain multiplicity one or be thickened
arbitrarily. The resulting bridgeless multigraph satisfies the one-sided
Merino--Welsh inequality whenever the stated no-bridge condition holds.

The connected counterexample (5.2) is graphic: it is a four-cycle with one
edge duplicated. Its heavy-edge support has three components, so the
two-component threshold in Corollary~\ref{cor:graphic} is also sharp under connectivity
alone.

\section{Computational checks}


The proof is symbolic. The accompanying exact Python verifier checks the
parallel-class activity formula, the corank-one induction interfaces, equality
controls, and both corank-two counterexamples on finite low-rank instances.
It uses explicit failures rather than Python assertions, so optimization mode
does not remove the checks.

From the published verification directory, run:

\begin{quote}\footnotesize\ttfamily
python\ \allowbreak{}-\allowbreak{}B\ \allowbreak{}verify\_\allowbreak{}corank\_\allowbreak{}one\_\allowbreak{}parallel\_\allowbreak{}support.py\par
\end{quote}

\begin{quote}\footnotesize\ttfamily
python\ \allowbreak{}-\allowbreak{}B\ \allowbreak{}-\allowbreak{}O\ \allowbreak{}verify\_\allowbreak{}corank\_\allowbreak{}one\_\allowbreak{}parallel\_\allowbreak{}support.py\par
\end{quote}


These computations test formulas and boundary examples; they are not a proof
of Theorem~\ref{thm:main}.

\section{Scope and prior work}


Noble and Royle proved the multiplicative Merino--Welsh inequality for
series--parallel graphs \cite{noble-royle}. Ferroni and Schröter proved it for split matroids
\cite{ferroni-schroeter}. The class in Theorem~\ref{thm:main} is not contained in either class: for example,
doubling every edge of \(K_5\) satisfies the spanning-support hypothesis, is
not series--parallel, and has \(M(K_5)\) as a minor; since split matroids are
minor closed and \(M(K_5)\) is not split, this example is not split.

Kung's density and cocircuit conditions \cite{kung} cover some dense or large-cocircuit
subfamilies but not the full sparse class considered here. Csikvári recently
improved a related Tutte-evaluation inequality used in this area \cite{csikvari}. The activity
and deletion--contraction methods are standard; the proposed contribution is the
parallel-support rank threshold, its strict corank-one form, the complete
equality classification, and the sharp rank-two-support boundary.

The theorem proves only the one-sided inequality \(T_M(0,2)\ge T_M(1,1)\)
on the stated class. It does not prove the additive or multiplicative
strengthenings. In fact, the uniformly doubled matroids used by Beke,
Csáji, Csikvári and Pituk \cite{beke-et-al} belong to the spanning-support layer and show
that the multiplicative inequality cannot be added to Theorem~\ref{thm:main} in general.
Nor does this paper prove the full graph conjecture or the general conjecture
for matroids with two disjoint bases.

Directed searches completed on 4 September 2026 found no source stating the
same parallel-support rank theorem. This is a bounded novelty assessment, not
a guarantee of global priority or a substitute for external peer review.

\section*{AI-assisted research disclosure}


Carptopus is the responsible author and takes responsibility for the claims,
proofs, source use, and released artifacts. OpenAI Codex was used for
AI-assisted literature organization, proof development and stress testing,
exact verification, adversarial auditing, and manuscript preparation.

\begin{thebibliography}{9}
\footnotesize
\raggedright

\bibitem{merino-welsh}
C.~Merino and D.~J.~A.~Welsh,
\emph{Forests, colourings and acyclic orientations of the square lattice},
Annals of Combinatorics \textbf{3} (1999), 417--429.
\href{https://doi.org/10.1007/s000260050019}{doi:10.1007/s000260050019}.

\bibitem{noble-royle}
S.~D.~Noble and G.~F.~Royle,
\emph{The Merino--Welsh conjecture holds for series--parallel graphs},
European Journal of Combinatorics \textbf{38} (2014), 24--35.
\href{https://arxiv.org/abs/1303.6416}{arXiv:1303.6416}.

\bibitem{ferroni-schroeter}
L.~Ferroni and B.~Schr\"oter,
\emph{The Merino--Welsh conjecture for split matroids},
Annals of Combinatorics \textbf{27} (2023), 737--748.
\href{https://arxiv.org/abs/2204.07132}{arXiv:2204.07132}.

\bibitem{beke-et-al}
C.~Beke, G.~K.~Cs\'aji, P.~Csikv\'ari, and S.~Pituk,
\emph{The Merino--Welsh conjecture is false for matroids},
Advances in Mathematics \textbf{446} (2024), 109674.
\href{https://arxiv.org/abs/2311.01932}{arXiv:2311.01932}.

\bibitem{chavez-lomeli-et-al}
L.~E.~Ch\'avez-Lomel\'i, C.~Merino, S.~D.~Noble, and M.~Ram\'irez-Ib\'a\~nez,
\emph{Some inequalities for the Tutte polynomial},
European Journal of Combinatorics \textbf{32} (2011), 422--433.
\href{https://arxiv.org/abs/1004.2639}{arXiv:1004.2639}.

\bibitem{kung}
J.~P.~S.~Kung,
\emph{Inconsequential results on the Merino--Welsh conjecture for Tutte polynomials},
Australasian Journal of Combinatorics \textbf{79} (2021), 12--15.
\href{https://arxiv.org/abs/2105.01825}{arXiv:2105.01825}.

\bibitem{csikvari}
P.~Csikv\'ari,
\emph{Around the Merino--Welsh conjecture: improving Jackson's inequality},
arXiv:2502.19196 (2025).
\href{https://arxiv.org/abs/2502.19196}{arXiv:2502.19196}.

\end{thebibliography}

\end{document}
